\documentclass{shirobyoushi}
% 査読用に著者情報を伏せるときは blind オプションを付ける:
%   \documentclass[blind]{shirobyoushi}

\usepackage{luatexja-fontspec}
% 欧文フォント
\setmainfont{Times New Roman}
\setsansfont[BoldFont={MS Gothic}]{MS Gothic}

% 和文フォント（明朝＝本文，ゴシック＝見出し）
\setmainjfont[BoldFont={MS Mincho}]{MS Mincho}
\setsansjfont[BoldFont={MS Gothic}]{MS Gothic}

\usepackage{multirow}
\usepackage{amsmath}
\usepackage{mathtools}
\mathtoolsset{showonlyrefs=false}
\usepackage{unicode-math}
\setmathfont{STIX Two Math}
\usepackage{siunitx}

%% ============================================================
%% タイトル情報
%% ============================================================
\jtitle{空気調和・衛生工学会論文集 版下原稿作成用 \\ レイアウト見本}
%% 連報の場合は，\jreport{第1報} のように記入する。連報でない場合は，\jreport{} として空にする。
\jreport{第◯報}
\jsubtitle{和文論文作成例}

\etitle{Template for Japanese Manuscript for Transactions of the Society
of Heating, Air-Conditioning and Sanitary Engineers of Japan}
\ereport{Part ○}
\esubtitle{Template for English Abstract}

%% 著者: \addauthor{和文氏名}{英文氏名}{所属番号}
\addauthor{空調 太郎}{Taro KUCHO}{1}
\addauthor{空調 花子}{Hanako KUCHO}{2}
\addauthor{衛生 小次郎}{Kojiro EISEI}{3}
\addauthor{衛生 学}{Manabu EISEI}{3}

%% 所属: \addaffiliation{番号}{和文所属}{和文資格}{英文所属}{英文資格}
\addaffiliation{1}{空調工業（株）設備部}{特別会員}%
{Engineering Division, Kucho Industry Co., Ltd.}{Fellow Member}
\addaffiliation{2}{空調大学工学部}{正会員}%
{Faculty of Engineering, Kucho University}{Member}
\addaffiliation{3}{衛生大学大学院工学研究科}{SHASE 技術フェロー}%
{Graduate School of Engineering, Eisei University}{Fellow Engineer}

%% 和文概要（明朝9pt・37字詰・8行以内・300字程度）
\jabstract{
  論文委員会にて採用が決定した論文は，採用をお知らせするとともに，指定する期日までに，原則として完全な版下原稿を提出していただきます。この“レイアウト見本”は，当学会論文集の完全版下原稿（和文）を作成するために必要な情報を，版下原稿の体裁をとって記載してあります。論文投稿時および採用決定後は，“執筆要綱”とこの“レイアウト見本”に基づき，印刷仕上がり後に十分判読できるよう，版下原稿を精密かつ鮮明に作成してください。なお，“論文題名”，“著者名”，“勤務先・所属部署名および会員資格”につきましては，誌面統一のために，筆者作成内容に基づいて当学会にて組版し直しますが，スペースを正しくとり，正確に作成してください。
}

%% 英文概要（Synopsis，ローマン9pt・2段・550語以内）
\eabstract{
  最終ページに，和文論文の場合は英文概要を記入します。
  本文の最後のページの余白に概要の全文が収まる場合は，参考文献の下に約12mmの間隔を空けて続けて記入してください。
  全文が収まらない場合は，最終ページの先頭から1行分空間を空けて英文概要を記入します。
  英文概要は本文と同じスペースに2段組みで550語以内（1ページ以内）とし，“Synopsis:”で始めます。
  文字の種別はローマン体・9ポイントとします。
  文末に，事務局より通知する原稿受理日を下の例のように記入します。
}

%% キーワード（和文は中黒区切り，英文はカンマ区切り。手法KWを先頭に置く）
\jkeywords{理論解析・実験・実測}
\ekeywords{Theoretical Analysis, Experiment, Measurement}

%% 受理日（事務局が記入）
\ereceiveddate{Received October 30, 1998}

\begin{document}
\maketitle

%% ============================================================
\section*{はじめに}

本文は，明朝体・9ポイントで，1ページ26文字×46行の二段組みとします。

見出しのレベルは次のとおりとします。各段落の最初の行は，1文字分スペースを空けてください。

なお，“はじめに” ，“結論” ，“謝辞”および“付録”などは章番号を振りません。

%% ============================================================
\section{章の見出し（2行になるときには，この例のように改行します）}

\subsection{節の見出し}

\subsubsection{項の見出し}

\paragraph{細目の見出し}

以下，文章が続きます。
なお，章の見出しのみ上に1行空けてください。

箇条書きにする場合は，下記のとおりとします。
\begin{enumerate}
  \item 第1箇条書き
    \begin{enumerate}
      \item 第2箇条書き
    \end{enumerate}
\end{enumerate}

%% ============================================================
\section{数式}

文中の式については執筆要綱“12．数と数式”を参照して記入してください。

独立した式は数式用のフォントを用い，左から2文字分スペースを空けて，整数式は1行分，分数式は2行分のスペースの中央に記入し，右詰めで
$\cdots\cdots(1)$，$\cdots\cdots(2)$，のように，2文字分のリーダー
($\cdots$) に続けて番号を振ります。

量記号で構成された式の掛け算は“$\cdot$”や“$\times$”を用いずに記入し，主に数値で構成された式の掛け算は“$\times$”を用いて記入します。

なお，2行以上になる数式は，数学記号が行の最初になるように工夫して改行します。
%
\begin{equation}
  Q_0 \frac{\mathrm{d}\theta_r}{\mathrm{d}t}
  = W_0 (\theta_\mathrm{out}-\theta_r^\prime)+\varepsilon S + H
  \label{eq:heat}
\end{equation}
%
ここに，
\begin{itemize}
  \item[] $W_0$: 基準システムの総括熱通過量 \hfill (\unit{\watt\per\kelvin})
  \item[] $Q_0$: 基準システムの熱容量 \hfill (\unit{\watt})
  \item[] $S$: 日射量 \hfill (\unit{\watt})
  \item[] $H$: 吸収熱量 \hfill (\unit{\watt})
  \item[] $\varepsilon$: 集熱効率 \hfill (-)
  \item[] $\theta_r^\prime$: 基準システムの基準計算室温 \hfill (\unit{\celsius})
  \item[] $\theta_\mathrm{out}$: 測定外気温 \hfill (\unit{\celsius})
\end{itemize}

量記号は，執筆要綱 “11. 記号の書き方” のとおり用い，英語の頭文字を用いたりしません。

単位は，右詰めで記入します。

%% ============================================================
\section{図・写真および表}

図・写真および表は執筆要綱 “10．図・写真および表”
とこのレイアウト見本を参照して作成してください\quotenote{脚注は，本文中の注をつける部分に*を付し，当該ページおよび当該の片段の最終行の下に罫線を引き，この作成例のように記入します。文字は，原則的に明朝体・8ポイントを使用してください。当該ページに複数の脚注がある場合は*1，*2，…としてください。}。
なお，タイトルの位置は，図・写真は図もしくは写真の下側中央とし，表は表の上側中央とします。文字の大きさは8ポイントとします。また，図・写真の注や凡例は図・写真のタイトルと図・写真の間に，表の注や凡例は表の下に8ポイントで記入する。

図・写真および表を配置するのは，論文中で最初に引用されている当該ページ，またはその前後2ページ以内とします。原則としてページの上方にバランス良く配置してください。本文との間は，1行空けます。

大きさは，片段または通し（両段）の二種類とし，最大は1ページ（所定原稿用紙の枠内）大とする。横幅を片段または両段分より狭くして，図・写真および表の横に文章を配置してはいけません。

図および表を本文中に引用する際は，ゴシック体で，\textgt{表-\ref{tab:sample}}，\textgt{図-\ref{fig:sample1}}のように記入します。

% 図表は原則として片段（単段）に収める。
% 両段に収める場合は table* 環境を用いる。
\begin{table*}[tb]
  \centering
  \caption{表のタイトルは表の上側中央に配置する}\label{tab:sample}
  \small
  \setlength{\tabcolsep}{6pt}
  \begin{tabular}{
      c l c
      S[table-format=-3.2]
      S[table-format=-3.2]
      S[table-format=-3.2]
      S[table-format=-2.1]
    }
    \toprule
    & & ケース & {投入熱量 (\unit{\watt})} & {対流熱伝達量 (\unit{\watt})}
    & {放射熱伝達量\textsuperscript{\dag2} (\unit{\watt})} & {一様拡散温度
    (\unit{\celsius})} \\
    \midrule
    \multirow{4}{*}{外乱・内乱}
    & 窓面       & 1 & 360.33 & 759.44\rlap{\textsuperscript{\dag1}} &
    179.10 & 6.5 \\
    & 外壁       & 2 & 34.89  & 10.47  & 24.42  & 0.4 \\
    & 内部発熱体 & 3 & 116.30 & 116.30 & 0      & 4.2 \\
    & 全断熱壁   & 4 & 0      & 88.74  & -83.74 & 3.0 \\
    \midrule
    \multirow{3}{*}{操作要因}
    & 天井冷放射面 & 5 & -232.60 & -120.95 & -111.65 & -4.3 \\
    & 吹出し口     & 6 & -279.12 & -279.12 & 0       & -10.0 \\
    & 吸込み口     & 7 & 0       & 0       & 0       & 0.0 \\
    \bottomrule
  \end{tabular}
  \tablenote{
    \textgt{注}\quad\textsuperscript{\dag1}
    表中の特定の数値や項目に注が必要な場合は，上付きの\dag（ダガー）を付し，表中と対応させて書きます。\par
    \hphantom{\textgt{注}\quad}\textsuperscript{\dag2}
    表中の数値のうち，関連のある数値はけた数，小数点をそろえてください。\par
  \hphantom{\textgt{注}\quad}表全体に対する注は，\dag（ダガー）をつける必要はありません。}
\end{table*}

\begin{figure}[htbp]
  \centering
  \includegraphics[width=\linewidth]{example-image}
  \figurenote{\textgt{注} 図の注や説明がある場合は，この位置に左詰めで記入します。}
  \caption{図のタイトルは図の下中央に配置する}\label{fig:sample1}
\end{figure}

%% ============================================================
\section*{結論}
結論がある場合は，このように“結論” の見出しをつけて章番号を振らずに組み入れます。

%% ============================================================
\acknowledgments%
謝辞がある場合は“結論” の後に組み入れます。
章番号は振りません。

%% ============================================================
\appendix%
付録がある場合は，“結論”，“謝辞” の後に組み入れます。
章番号は振りません。

%% ============================================================
%% 参考文献（執筆要領 13.5 の体裁）
%%   references.bib に BibTeX 形式で記述し，upbibtex + junsrt.bst で組版。
%%   \nocite{*} で .bib 記載順に全件を出力する。
%% ============================================================
\nocite{*}
\bibliographystyle{junsrt}
\bibliography{references}

%% ============================================================
%% 英文要素（最終ページ）。\end{document} の直前に置く。
%% ============================================================
\makeenglishabstract%

\end{document}
